\documentclass[10pt]{article}
\usepackage[a4paper,margin=0.55in]{geometry}
\usepackage{amsmath,amssymb}
\usepackage{xcolor}
\usepackage{hyperref}
\setlength{\jot}{5pt}
\numberwithin{equation}{section}
\allowdisplaybreaks

\title{UV divergence of $N_{2}$ (Group I, Row III) following Balitsky--Chirilli and Lublinsky--Mulian:\\ subtraction in the Wilson lines only, dimensional regularisation of the subtraction term}
\date{}
\begin{document}
\maketitle

\section{Starting expression}
\begin{align}
    \frac{d^{3}N_{2\,\rm NLO}^{\rm}}{d^{3}k}  & =  \frac{1}{(2\pi)^{3}}\, \frac{g^{4}}{8\pi^{4}}\,\frac{1}{k^{+}}\,\int_{\Lambda}^{k^{+}-\Lambda}\, \frac{dp^{+}}{2\pi} \, \int_{x,y,z,x',w'}\,e^{-ik\cdot(w'-z)}\,e^{-ik\cdot (z-x)\,\frac{p^{+}}{k^{+}}}\,\frac{1}{\left[p^{+}\,(y-x)^{2}+(k^{+}-p^{+})\,(y-z)^{2}\right]}\notag \\ &\times \frac{(x'-w')^{i}\,(z-x)^{m}}{(x'-w')^{2}\,(z-x)^{2}}\Bigg[\frac{\delta^{im}\,d_{\perp}}{2\,(x-z)^{2}}\,\left[(y-x)^{2}-(y-z)^{2}\right]\,\frac{p^{+}\,(k^{+}-p^{+})}{k^{+\,2}}\notag\\ &-\frac{p^{+}}{(k^{+}-p^{+})}\,\left[\frac{(y-x)^{i}\,(x-z)^{m}}{(x-z)^{2}}-\frac{(y-x)^{i}\,(y-z)^{m}}{2\,(y-z)^{2}}\right]\notag \\ & +\frac{(k^{+}-p^{+})\,\delta_{im}}{k^{+}}\,\left[\frac{(y-z)\cdot (x-z)}{(x-z)^{2}}+\frac{(y-x)\cdot (y-z)}{2\,(y-x)^{2}}-\frac{\left[(y-x)^{2}-(y-z)^{2}\right]}{2\,(x-z)^{2}}\right]\notag \\ & -\frac{(k^{+}-p^{+})}{p^{+}}\,\left[\frac{(y-z)^{i}\,(x-z)^{m}}{(x-z)^{2}}+\frac{(y-x)^{m}\,(y-z)^{i}}{2\,(y-x)^{2}}\right]\notag \\ & +\frac{p^{+}\,\delta_{im}}{k^{+}}\,\left[\frac{(y-x)\cdot (x-z)}{(x-z)^{2}}-\frac{(y-x)\cdot (y-z)}{2\,(y-z)^{2}}-\frac{\left[(y-x)^{2}-(y-z)^{2}\right]}{2\,(x-z)^{2}}\right]\notag \\ &-\left[\frac{(y-x)^{m}\,(x-z)^{i}}{(x-z)^{2}}-\frac{(y-x)^{m}\,(y-z)^{i}}{2\,(y-z)^{2}}+\frac{(y-z)^{m}\,(x-z)^{i}}{(x-z)^{2}}+\frac{(y-x)^{i}\,(y-z)^{m}}{2\,(y-x)^{2}}\right]\Bigg]\notag \\ & \times\Big[U^{ab'}(x')-U^{ab'}(w')\Big]\notag\\&\times\Big[-f^{adc}\,f^{fbe}\,U^{db}(x)\,U^{ce}(z)+N_{c}\,U^{af}(y)+N_{c}\,U^{af}(x) -N_{c}\,U^{af}(x)\Big]\,\rho^{b'}(x')\,\rho^{f}(y).
\label{eq:N2}
\end{align}
From here on we use $\xi=p^{+}/k^{+}$, so
\[\int_{\Lambda}^{k^{+}-\Lambda}\frac{dp^{+}}{2\pi}\frac{1}{p^{+}(y-x)^{2}+(k^{+}-p^{+})(y-z)^{2}}=\int_{\Lambda/k^{+}}^{1-\Lambda/k^{+}}\frac{d\xi}{2\pi}\frac{1}{\xi(y-x)^{2}+(1-\xi)(y-z)^{2}}.\]

\section{Step 1: subtraction in the Wilson lines only}
(Balitsky--Chirilli eq.~(59); Lublinsky--Mulian eqs.~(4.71)--(4.73).)
The only UV divergence is at $z\to x$, where the kernel goes as $1/(x-z)^{2}$. Following Balitsky--Chirilli and Lublinsky--Mulian, we subtract and add the configuration in which the two gluons cross the shockwave at the \emph{same} point. Only the Wilson lines are changed; the kernel, \emph{including the exponential}, is untouched. Setting $z\to x$ in the colour factor and using
\begin{equation}
f^{abc}\,f^{def}\,U^{be}(x)\,U^{cf}(x)=N_{c}\,U^{ad}(x)
\end{equation}
gives $-f^{adc}f^{fbe}U^{db}(x)U^{ce}(x)+N_{c}U^{af}(y)=N_{c}\big[U^{af}(y)-U^{af}(x)\big]$, the single-gluon colour structure. Hence
\begin{equation}
\frac{d^{3}N_{2}}{d^{3}k}=\underbrace{\Big(\frac{d^{3}N_{2}}{d^{3}k}-\frac{d^{3}N^{\rm sub}_{2}}{d^{3}k}\Big)}_{[-ffU(x)U(z)+N_{c}U(x)]}+\underbrace{\frac{d^{3}N^{\rm sub}_{2}}{d^{3}k}}_{N_{c}[U(y)-U(x)]},
\end{equation}
with the same kernel in both terms. This is exactly the split $N_{2}=N_{2a}+N_{2b}$.
\begin{itemize}
\item $N_{2}-N_{2}^{\rm sub}$ is UV finite, because its colour factor vanishes at $z=x$. It is computed in $d=2$.
\item $N_{2}^{\rm sub}$ carries the UV divergence. Its colour factor does not depend on $z$, so the $z$ integral of the \emph{full kernel} is done in $d_{\perp}=2-\epsilon$.
\end{itemize}

\section{Step 2: exact split of the subtraction kernel}
Contract with $(x'-w')^{i}$. The bracket of \eqref{eq:N2}, multiplied by $\frac{(z-x)^{m}}{(z-x)^{2}}$, is identically equal to (UV part) $+$ (rest):
\begin{align}
\text{UV:}\quad&d_{\perp}\,\xi(1-\xi)\,\frac{(z-x)^{i}\,(z-x)\cdot(y-x)}{(z-x)^{4}}+\Big(\frac{\xi}{1-\xi}+\frac{1-\xi}{\xi}\Big)\frac{(y-x)^{i}}{(z-x)^{2}},\\
\text{rest:}\quad&\frac{(z-x)^{i}}{(z-x)^{2}}\Big(-\frac{d_{\perp}}{2}\xi(1-\xi)+2(1-\xi)+\frac{\xi}{2}-(1-\xi)\frac{(y-x)\cdot(z-x)}{2(y-x)^{2}}-\xi\frac{(y-x)\cdot(y-z)}{2(y-z)^{2}}\Big)\notag\\
&+\frac{\xi}{1-\xi}\frac{(y-x)^{i}\,(z-x)\cdot(y-z)}{2(z-x)^{2}(y-z)^{2}}-\frac{1-\xi}{\xi}\Big(\frac{(z-x)^{i}}{(z-x)^{2}}+\frac{(z-x)\cdot(y-x)\,(y-z)^{i}}{2(z-x)^{2}(y-x)^{2}}\Big)\notag\\
&-\frac{(z-x)^{i}}{(z-x)^{2}}+\frac{(z-x)\cdot(y-x)\,(y-z)^{i}}{2(z-x)^{2}(y-z)^{2}}-\frac{(y-x)^{i}\,(z-x)\cdot(y-z)}{2(z-x)^{2}(y-x)^{2}} .
\end{align}
This was checked numerically at random points for $d_{\perp}=2.0,\,2.3,\,2.6,\,2.9$; the residual is $10^{-15}$. The rest is $O(1/|z-x|)$, so it is UV finite and taken in $d=2$. The UV part contains the whole UV divergence. Its angle-averaged coefficient is $\frac{d_{\perp}}{d}\xi(1-\xi)+\frac{\xi}{1-\xi}+\frac{1-\xi}{\xi}$, which equals $\xi(1-\xi)+\frac{\xi}{1-\xi}+\frac{1-\xi}{\xi}$ for $d_{\perp}=d$. We keep $d_{\perp}=d$ in the UV part.

\section{Step 3: dimensional regularisation of the UV part; the exponential becomes a momentum shift}
\textbf{3a. Energy denominator in momentum space.} By the four-dimensional identity $\int d^{4}P\,e^{iP\cdot R}/P^{2}=4\pi^{2}/R^{2}$, after rescaling $p\to\sqrt\xi\,p$, $l\to\sqrt{1-\xi}\,l$:
\begin{equation}
\frac{1}{\xi(y-x)^{2}+(1-\xi)(y-z)^{2}}=\frac{1}{4\pi^{2}\,\xi(1-\xi)}\int d^{2}p\,d^{2}l\;\frac{e^{\,ip\cdot(y-x)+il\cdot(y-z)}}{\frac{p^{2}}{\xi}+\frac{l^{2}}{1-\xi}} .
\end{equation}

\textbf{3b. The exponential.} Since $e^{-ik\cdot(w'-z)}e^{-ik\cdot(z-x)\xi}=e^{-ik\cdot(w'-x)}\,e^{i(1-\xi)k\cdot(z-x)}$ and $e^{il\cdot(y-z)}=e^{il\cdot(y-x)}e^{-il\cdot(z-x)}$, the whole $z$ dependence of the UV part is
\begin{equation}
\int d^{d}z\;e^{-i\,(l-(1-\xi)k)\cdot(z-x)}\Big[d_{\perp}\xi(1-\xi)\frac{(z-x)^{i}(z-x)^{j}}{(z-x)^{4}}(y-x)^{j}+\Big(\frac{\xi}{1-\xi}+\frac{1-\xi}{\xi}\Big)\frac{(y-x)^{i}}{(z-x)^{2}}\Big].
\end{equation}
The exponential is therefore not reduced or approximated. It just shifts the momentum at which the kernel is Fourier transformed, to $l-(1-\xi)k$. In Balitsky--Chirilli's dipole case there is no external momentum, so there is no shift.

\textbf{3c. The $d$-dimensional Fourier transforms} ($d=2-\epsilon$, $\overline{MS}$, $\bar\mu^{2}=4\pi e^{-\gamma_{E}}\mu^{2}$):
\begin{align}
\mu^{-\epsilon}\!\int d^{d}r\,\frac{e^{-iQ\cdot r}}{r^{2}}&=\pi^{d/2}\,2^{d-2}\,\Gamma\big(\tfrac d2-1\big)\,(Q^{2})^{1-\frac d2}\mu^{-\epsilon}=\pi\Big[\frac{2}{\epsilon}+\log\frac{\bar\mu^{2}}{Q^{2}}\Big],\\
\mu^{-\epsilon}\!\int d^{d}r\,e^{-iQ\cdot r}\,\frac{r^{i}r^{j}}{r^{4}}&=\frac{\delta^{ij}}{d}\,\pi\Big[\frac{2}{\epsilon}+\log\frac{\bar\mu^{2}}{Q^{2}}\Big]-\pi\Big(\frac{Q^{i}Q^{j}}{Q^{2}}-\frac{\delta^{ij}}{2}\Big).
\end{align}
The traceless part follows from $\int d\varphi\,e^{-iQr\cos(\varphi-\varphi_{Q})}\cos2\varphi=-2\pi J_{2}(Qr)\cos2\varphi_{Q}$ and $\int_{0}^{\infty}J_{2}(x)\,dx/x=\frac12$. As a check, with a short-distance cut-off $\delta$, $\int d^{2}r\,e^{-iQ\cdot r}/r^{2}=2\pi[-\log(|Q|\delta/2)-\gamma_{E}]$, and the replacement $\log\frac{1}{\delta^{2}}\to\frac{2}{\epsilon}+2\gamma_{E}+\log\frac{\bar\mu^{2}}{4}$ reproduces the first line.

\textbf{3d. Result for the UV part} ($d_{\perp}=d$):
\begin{align}
\frac{d^{3}N^{\rm sub}_{2\,\rm UV}}{d^{3}k}&=\frac{1}{(2\pi)^{3}}\,\frac{g^{4}}{8\pi^{4}}\,\frac{1}{k^{+}}\int_{\Lambda/k^{+}}^{1-\Lambda/k^{+}}\frac{d\xi}{2\pi}\int_{x,y,x',w'}e^{-ik\cdot(w'-x)}\notag\\&\times\frac{1}{4\pi^{2}\,\xi(1-\xi)}\int d^{2}p\,d^{2}l\;\frac{e^{\,i(p+l)\cdot(y-x)}}{\frac{p^{2}}{\xi}+\frac{l^{2}}{1-\xi}}\;\frac{1}{(x'-w')^{2}}\notag\\
&\times\Bigg\{\pi\Big(\xi(1-\xi)+\frac{\xi}{1-\xi}+\frac{1-\xi}{\xi}\Big)\,(x'-w')\cdot(y-x)\left[\frac{2}{\epsilon}+\log\frac{\bar\mu^{2}}{\big(l-(1-\xi)k\big)^{2}}\right]\notag\\
&\qquad-2\pi\,\xi(1-\xi)\left[\frac{(x'-w')\cdot\big(l-(1-\xi)k\big)\;\big(l-(1-\xi)k\big)\cdot(y-x)}{\big(l-(1-\xi)k\big)^{2}}-\frac{(x'-w')\cdot(y-x)}{2}\right]\Bigg\}\notag\\
&\times\Big[U^{ab'}(x')-U^{ab'}(w')\Big]\,N_{c}\Big[U^{af}(y)-U^{af}(x)\Big]\rho^{b'}(x')\,\rho^{f}(y).
\label{eq:UVmom}
\end{align}
If desired, the $p$ integral can be done: $\int d^{2}p\,\frac{e^{ip\cdot(y-x)}}{p^{2}/\xi+l^{2}/(1-\xi)}=2\pi\xi\,K_{0}\big(|l|\,|y-x|\sqrt{\xi/(1-\xi)}\big)$.

The equivalent coordinate-space form (with the $l$ integral done using a Feynman parameter and Bessel functions $K_{0},K_{1},K_{2}$) is in \texttt{N2b\_UV\_z\_integrated\_explicit.pdf}. The two forms were checked against each other numerically; they agree to the accuracy of the momentum grid ($2\times10^{-3}$).

\textbf{3e. The UV pole and the $\log\bar\mu$.} Undoing the $p,l$ integral at fixed scale gives $1/(y-x)^{2}$. So the $\frac{2}{\epsilon}$ and the $\log\bar\mu^{2}$ multiply
\[
\pi\Big(\xi(1-\xi)+\frac{\xi}{1-\xi}+\frac{1-\xi}{\xi}\Big)\frac{(x'-w')\cdot(y-x)}{(x'-w')^{2}(y-x)^{2}} .
\]
With $\int_{\Lambda/k^{+}}^{1-\Lambda/k^{+}}d\xi\,\big(\xi(1-\xi)+\frac{\xi}{1-\xi}+\frac{1-\xi}{\xi}\big)=2\log\frac{k^{+}-\Lambda}{\Lambda}-\frac{11}{6}+\frac{4\Lambda}{k^{+}}-\frac{\Lambda^{2}}{k^{+2}}+\frac{2\Lambda^{3}}{3k^{+3}}$, the coefficient of $\log\bar\mu^{2}$ is
\begin{align}
&\frac{1}{(2\pi)^{3}}\,\frac{g^{4}N_{c}}{16\pi^{4}}\,\frac{1}{k^{+}}\Big[2\log\frac{k^{+}}{\Lambda}-\frac{11}{6}\Big]\int_{x,y,x',w'}e^{-ik\cdot(w'-x)}\frac{(x'-w')\cdot(y-x)}{(x'-w')^{2}(y-x)^{2}}\notag\\&\qquad\times\Big[U^{ab'}(x')-U^{ab'}(w')\Big]\Big[U^{af}(y)-U^{af}(x)\Big]\rho^{b'}(x')\rho^{f}(y).
\end{align}
In the convention $\int d^{2}p/p^{2}=0$ the $\frac{2}{\epsilon}$ is dropped. The $-\frac{11}{6}\log\bar\mu^{2}$ is the running-coupling term, which combines with Group~I Row~II. The $2\log\frac{k^{+}}{\Lambda}\log\bar\mu^{2}$ goes with the JIMWLK terms.

\section{Step 4: the finite pieces}
\textbf{4a. $N_{2}-N_{2}^{\rm sub}$} ($d=2$, $d_{\perp}=2$): the full expression \eqref{eq:N2} with colour factor
\[
\Big[U^{ab'}(x')-U^{ab'}(w')\Big]\Big[-f^{adc}f^{fbe}U^{db}(x)U^{ce}(z)+N_{c}U^{af}(x)\Big]\rho^{b'}(x')\rho^{f}(y).
\]

\textbf{4b. Rest of the subtraction term} ($d=2$, $d_{\perp}=2$):
\begin{align}
\frac{d^{3}N^{\rm sub}_{2\,\rm rest}}{d^{3}k}&=\frac{1}{(2\pi)^{3}}\,\frac{g^{4}}{8\pi^{4}}\,\frac{1}{k^{+}}\int_{\Lambda/k^{+}}^{1-\Lambda/k^{+}}\frac{d\xi}{2\pi}\int_{x,y,z,x',w'}e^{-ik\cdot(w'-z)}\,e^{-ik\cdot(z-x)\xi}\notag\\&\times\frac{1}{\xi(y-x)^{2}+(1-\xi)(y-z)^{2}}\,\frac{(x'-w')^{i}}{(x'-w')^{2}}\notag\\&\times\Bigg[\frac{(z-x)^{i}}{(z-x)^{2}}\Big(-\xi(1-\xi)+2(1-\xi)+\frac{\xi}{2}-(1-\xi)\frac{(y-x)\cdot(z-x)}{2(y-x)^{2}}-\xi\frac{(y-x)\cdot(y-z)}{2(y-z)^{2}}\Big)\notag\\
&\quad+\frac{\xi}{1-\xi}\frac{(y-x)^{i}\,(z-x)\cdot(y-z)}{2(z-x)^{2}(y-z)^{2}}-\frac{1-\xi}{\xi}\Big(\frac{(z-x)^{i}}{(z-x)^{2}}+\frac{(z-x)\cdot(y-x)\,(y-z)^{i}}{2(z-x)^{2}(y-x)^{2}}\Big)\notag\\
&\quad-\frac{(z-x)^{i}}{(z-x)^{2}}+\frac{(z-x)\cdot(y-x)\,(y-z)^{i}}{2(z-x)^{2}(y-z)^{2}}-\frac{(y-x)^{i}\,(z-x)\cdot(y-z)}{2(z-x)^{2}(y-x)^{2}}\Bigg]\notag\\
&\times\Big[U^{ab'}(x')-U^{ab'}(w')\Big]\,N_{c}\Big[U^{af}(y)-U^{af}(x)\Big]\rho^{b'}(x')\,\rho^{f}(y).
\end{align}

\section{Summary}
\begin{equation}
\frac{d^{3}N_{2\,\rm NLO}}{d^{3}k}=\underbrace{\Big(\frac{d^{3}N_{2}}{d^{3}k}-\frac{d^{3}N^{\rm sub}_{2}}{d^{3}k}\Big)}_{\text{finite, }d=2}+\underbrace{\frac{d^{3}N^{\rm sub}_{2\,\rm UV}}{d^{3}k}}_{\text{eq. \eqref{eq:UVmom}, dim.\ reg.}}+\underbrace{\frac{d^{3}N^{\rm sub}_{2\,\rm rest}}{d^{3}k}}_{\text{finite, }d=2}.
\end{equation}
\begin{itemize}
\item As in BC and LM, only the Wilson lines are changed by the subtraction; the exponential stays in the kernel. In the $d$-dimensional $z$ integral it becomes the momentum shift $l\to l-(1-\xi)k$. No extra subtraction in $z$ is made.
\item The only $\log\bar\mu$ is in eq.~\eqref{eq:UVmom}, with coefficient $\frac{g^{4}N_{c}}{16\pi^{4}}\big[2\log\frac{k^{+}}{\Lambda}-\frac{11}{6}\big]$.
\item \textbf{Next step: the plus prescription in $\xi$} on the finite pieces and on the remaining $\xi$ dependence of eq.~\eqref{eq:UVmom}.
\item \textbf{Caution 1.} In the rest, the term $-\frac{1-\xi}{\xi}\frac{(z-x)^{i}}{(z-x)^{2}}$ gives a spurious divergence at $z\to y$ in the $\xi\to0$ endpoint, where the denominator becomes $(y-z)^{2}$. It cancels against the UV part. To avoid it, move this term into the UV part, where it turns $\frac{1-\xi}{\xi}\frac{(y-x)^{i}}{(z-x)^{2}}$ into $\frac{1-\xi}{\xi}\frac{(y-z)^{i}}{(z-x)^{2}}$. The UV properties of both parts are unchanged.
\item \textbf{Caution 2.} Near $\xi\to1$ the momentum shift vanishes, $(1-\xi)k\to0$, and the $\xi\to1$ endpoint has large-distance (infrared) behaviour. This shows up as $\frac12\log^{2}\frac{k^{+}}{\Lambda}$ in the coordinate-space form. It must cancel against the other $\log(k^{+}/\Lambda)$ (JIMWLK) terms; it is not a UV effect.
\end{itemize}

\end{document}
